% !PRINTED 397
\weberpage{397}
%<<CONT
sein, weil 2) als reeller Complex keine reelle Doppeltangente
mit einer imaginären gepaart enthalten kann.

Es sei also
\[
2)\quad \xi\xi',\ x_1x_2,\ y_1y_2,
\]
und daraus leiten wir die zwei Complexe her:
\[
\begin{aligned}
3)&\quad x_1\xi,\ x_2\xi'\\
4)&\quad x_1\xi',\ x_2\xi,
\end{aligned}
\]
die mit einander syzygetisch sind, und mithin ausser den vier
angegebenen kein Element gemein haben. Sie sind aber zugleich
conjugirt, und daher kann in 3) ausser $x_1, x_2$ keine reelle Doppeltangente
 vorkommen. Nun sind aber 1) und 3) azygetisch, weil
$\xi$ in 1) nicht vorkommt; mithin muss aus jedem Paare von 1)
ein Element in 3) vorkommen. Also enthält 3) ausser $x_1, x_2$ noch
reelle Elemente, wodurch sich der Widerspruch ergiebt.

Wir schliessen daraus, dass die beiden Complexe, die durch
die Paare $x_1x_2$, $y_1y_2$ bestimmt sind und mit 1) ein syzygetisches
Complextripel ausmachen, \gesperrt{nur reelle Elemente} enthalten.

Lassen wir einen dieser Complexe, etwa $x_1x_2$, an Stelle
von 1) treten und wiederholen dann unseren Schluss, so ergiebt
sich, dass überhaupt alle Doppeltangenten, also auch die des
letzten Paares von 1), reell sein müssen.

5. \gesperrt{Wenn mehr als acht reelle Doppeltangenten
vorhanden sind, so giebt es sechzehn, die in
einem syzygetischen Complextripel je vier reelle
Paare bilden.}

Nehmen wir mehr als acht reelle Doppeltangenten an, so
können wir nach den Sätzen 3., 4. folgendes syzygetische Complextripel
 zusammenstellen:
\[
\begin{aligned}
1)&\quad x_1y_1,\ x_2y_2,\ x_3y_3,\ x_4y_4\\
2)&\quad x_1x_2,\ y_1y_2,\ z_1z_2\\
3)&\quad x_1y_2,\ x_2y_1,
\end{aligned}
\]
worin die $x_i, y_i, z_i$ reell sind.

Damit verbinden wir den Complex
\[
4)\quad x_1z_1,\ x_2z_2,
\]
der mit 1) azygetisch ist, und daher aus den Paaren $x_3y_3, x_4y_4$
je ein und nur ein Element enthalten kann, etwa $x_3$ und $x_4$.
Da aber 4) ein reeller Complex ist, so muss es zwei weitere
%CONT>>
